\chapter{Web Interface}

% === Source file: web/control-ui.md ===
\section{Control UI (browser)}


The Control UI is a small \textbf{Vite + Lit} single-page app served by the Gateway:

\begin{itemize}
  \item default: \texttt{http://:18789/}
  \item optional prefix: set \texttt{gateway.controlUi.basePath} (e.g. \texttt{/openclaw})
\end{itemize}


It speaks \textbf{directly to the Gateway WebSocket} on the same port.

\subsection{Quick open (local)}


If the Gateway is running on the same computer, open:

\begin{itemize}
  \item \href{http://127.0.0.1:18789/}{http://127.0.0.1:18789/} (or \href{http://localhost:18789/}{http://localhost:18789/})
\end{itemize}


If the page fails to load, start the Gateway first: \texttt{openclaw gateway}.

Auth is supplied during the WebSocket handshake via:

\begin{itemize}
  \item \texttt{connect.params.auth.token}
  \item \texttt{connect.params.auth.password}
\end{itemize}

The dashboard settings panel lets you store a token; passwords are not persisted.
The onboarding wizard generates a gateway token by default, so paste it here on first connect.

\subsection{Device pairing (first connection)}


When you connect to the Control UI from a new browser or device, the Gateway
requires a \textbf{one-time pairing approval} — even if you're on the same Tailnet
with \texttt{gateway.auth.allowTailscale: true}. This is a security measure to prevent
unauthorized access.

\textbf{What you'll see:} "disconnected (1008): pairing required"

\textbf{To approve the device:}

\begin{lstlisting}[language=bash]
# List pending requests
openclaw devices list

# Approve by request ID
openclaw devices approve <requestId>
\end{lstlisting}


Once approved, the device is remembered and won't require re-approval unless
you revoke it with \texttt{openclaw devices revoke --device  --role }. See
\href{/cli/devices}{Devices CLI} for token rotation and revocation.

\textbf{Notes:}

\begin{itemize}
  \item Local connections (\texttt{127.0.0.1}) are auto-approved.
  \item Remote connections (LAN, Tailnet, etc.) require explicit approval.
  \item Each browser profile generates a unique device ID, so switching browsers or
\end{itemize}

clearing browser data will require re-pairing.

\subsection{What it can do (today)}


\begin{itemize}
  \item Chat with the model via Gateway WS (\texttt{chat.history}, \texttt{chat.send}, \texttt{chat.abort}, \texttt{chat.inject})
  \item Stream tool calls + live tool output cards in Chat (agent events)
  \item Channels: WhatsApp/Telegram/Discord/Slack + plugin channels (Mattermost, etc.) status + QR login + per-channel config (\texttt{channels.status}, \texttt{web.login.*}, \texttt{config.patch})
  \item Instances: presence list + refresh (\texttt{system-presence})
  \item Sessions: list + per-session thinking/verbose overrides (\texttt{sessions.list}, \texttt{sessions.patch})
  \item Cron jobs: list/add/edit/run/enable/disable + run history (\texttt{cron.*})
  \item Skills: status, enable/disable, install, API key updates (\texttt{skills.*})
  \item Nodes: list + caps (\texttt{node.list})
  \item Exec approvals: edit gateway or node allowlists + ask policy for \texttt{exec host=gateway/node} (\texttt{exec.approvals.*})
  \item Config: view/edit \texttt{\textasciitilde{}/.openclaw/openclaw.json} (\texttt{config.get}, \texttt{config.set})
  \item Config: apply + restart with validation (\texttt{config.apply}) and wake the last active session
  \item Config writes include a base-hash guard to prevent clobbering concurrent edits
  \item Config schema + form rendering (\texttt{config.schema}, including plugin + channel schemas); Raw JSON editor remains available
  \item Debug: status/health/models snapshots + event log + manual RPC calls (\texttt{status}, \texttt{health}, \texttt{models.list})
  \item Logs: live tail of gateway file logs with filter/export (\texttt{logs.tail})
  \item Update: run a package/git update + restart (\texttt{update.run}) with a restart report
\end{itemize}


Cron jobs panel notes:

\begin{itemize}
  \item For isolated jobs, delivery defaults to announce summary. You can switch to none if you want internal-only runs.
  \item Channel/target fields appear when announce is selected.
  \item Webhook mode uses \texttt{delivery.mode = "webhook"} with \texttt{delivery.to} set to a valid HTTP(S) webhook URL.
  \item For main-session jobs, webhook and none delivery modes are available.
  \item Advanced edit controls include delete-after-run, clear agent override, cron exact/stagger options,
\end{itemize}

agent model/thinking overrides, and best-effort delivery toggles.
\begin{itemize}
  \item Form validation is inline with field-level errors; invalid values disable the save button until fixed.
  \item Set \texttt{cron.webhookToken} to send a dedicated bearer token, if omitted the webhook is sent without an auth header.
  \item Deprecated fallback: stored legacy jobs with \texttt{notify: true} can still use \texttt{cron.webhook} until migrated.
\end{itemize}


\subsection{Chat behavior}


\begin{itemize}
  \item \texttt{chat.send} is \textbf{non-blocking}: it acks immediately with \texttt{\{ runId, status: "started" \}} and the response streams via \texttt{chat} events.
  \item Re-sending with the same \texttt{idempotencyKey} returns \texttt{\{ status: "in\_flight" \}} while running, and \texttt{\{ status: "ok" \}} after completion.
  \item \texttt{chat.history} responses are size-bounded for UI safety. When transcript entries are too large, Gateway may truncate long text fields, omit heavy metadata blocks, and replace oversized messages with a placeholder (\texttt{[chat.history omitted: message too large]}).
  \item \texttt{chat.inject} appends an assistant note to the session transcript and broadcasts a \texttt{chat} event for UI-only updates (no agent run, no channel delivery).
  \item Stop:
  \item Click \textbf{Stop} (calls \texttt{chat.abort})
  \item Type \texttt{/stop} (or standalone abort phrases like \texttt{stop}, \texttt{stop action}, \texttt{stop run}, \texttt{stop openclaw}, \texttt{please stop}) to abort out-of-band
  \item \texttt{chat.abort} supports \texttt{\{ sessionKey \}} (no \texttt{runId}) to abort all active runs for that session
  \item Abort partial retention:
  \item When a run is aborted, partial assistant text can still be shown in the UI
  \item Gateway persists aborted partial assistant text into transcript history when buffered output exists
  \item Persisted entries include abort metadata so transcript consumers can tell abort partials from normal completion output
\end{itemize}


\subsection{Tailnet access (recommended)}


\subsubsection{Integrated Tailscale Serve (preferred)}


Keep the Gateway on loopback and let Tailscale Serve proxy it with HTTPS:

\begin{lstlisting}[language=bash]
openclaw gateway --tailscale serve
\end{lstlisting}


Open:

\begin{itemize}
  \item \texttt{https:///} (or your configured \texttt{gateway.controlUi.basePath})
\end{itemize}


By default, Control UI/WebSocket Serve requests can authenticate via Tailscale identity headers
(\texttt{tailscale-user-login}) when \texttt{gateway.auth.allowTailscale} is \texttt{true}. OpenClaw
verifies the identity by resolving the \texttt{x-forwarded-for} address with
\texttt{tailscale whois} and matching it to the header, and only accepts these when the
request hits loopback with Tailscale’s \texttt{x-forwarded-*} headers. Set
\texttt{gateway.auth.allowTailscale: false} (or force \texttt{gateway.auth.mode: "password"})
if you want to require a token/password even for Serve traffic.
Tokenless Serve auth assumes the gateway host is trusted. If untrusted local
code may run on that host, require token/password auth.

\subsubsection{Bind to tailnet + token}


\begin{lstlisting}[language=bash]
openclaw gateway --bind tailnet --token "$(openssl rand -hex 32)"
\end{lstlisting}


Then open:

\begin{itemize}
  \item \texttt{http://:18789/} (or your configured \texttt{gateway.controlUi.basePath})
\end{itemize}


Paste the token into the UI settings (sent as \texttt{connect.params.auth.token}).

\subsection{Insecure HTTP}


If you open the dashboard over plain HTTP (\texttt{http://} or \texttt{http://}),
the browser runs in a \textbf{non-secure context} and blocks WebCrypto. By default,
OpenClaw \textbf{blocks} Control UI connections without device identity.

\textbf{Recommended fix:} use HTTPS (Tailscale Serve) or open the UI locally:

\begin{itemize}
  \item \texttt{https:///} (Serve)
  \item \texttt{http://127.0.0.1:18789/} (on the gateway host)
\end{itemize}


\textbf{Insecure-auth toggle behavior:}

\begin{lstlisting}[language=json]
{
  gateway: {
    controlUi: { allowInsecureAuth: true },
    bind: "tailnet",
    auth: { mode: "token", token: "replace-me" },
  },
}
\end{lstlisting}


\texttt{allowInsecureAuth} does not bypass Control UI device identity or pairing checks.

\textbf{Break-glass only:}

\begin{lstlisting}[language=json]
{
  gateway: {
    controlUi: { dangerouslyDisableDeviceAuth: true },
    bind: "tailnet",
    auth: { mode: "token", token: "replace-me" },
  },
}
\end{lstlisting}


\texttt{dangerouslyDisableDeviceAuth} disables Control UI device identity checks and is a
severe security downgrade. Revert quickly after emergency use.

See \href{/gateway/tailscale}{Tailscale} for HTTPS setup guidance.

\subsection{Building the UI}


The Gateway serves static files from \texttt{dist/control-ui}. Build them with:

\begin{lstlisting}[language=bash]
pnpm ui:build # auto-installs UI deps on first run
\end{lstlisting}


Optional absolute base (when you want fixed asset URLs):

\begin{lstlisting}[language=bash]
OPENCLAW_CONTROL_UI_BASE_PATH=/openclaw/ pnpm ui:build
\end{lstlisting}


For local development (separate dev server):

\begin{lstlisting}[language=bash]
pnpm ui:dev # auto-installs UI deps on first run
\end{lstlisting}


Then point the UI at your Gateway WS URL (e.g. \texttt{ws://127.0.0.1:18789}).

\subsection{Debugging/testing: dev server + remote Gateway}


The Control UI is static files; the WebSocket target is configurable and can be
different from the HTTP origin. This is handy when you want the Vite dev server
locally but the Gateway runs elsewhere.

\begin{enumerate}
  \item Start the UI dev server: \texttt{pnpm ui:dev}
  \item Open a URL like:
\end{enumerate}


\begin{lstlisting}
http://localhost:5173/?gatewayUrl=ws://<gateway-host>:18789
\end{lstlisting}


Optional one-time auth (if needed):

\begin{lstlisting}
http://localhost:5173/?gatewayUrl=wss://<gateway-host>:18789&token=<gateway-token>
\end{lstlisting}


Notes:

\begin{itemize}
  \item \texttt{gatewayUrl} is stored in localStorage after load and removed from the URL.
  \item \texttt{token} is stored in localStorage; \texttt{password} is kept in memory only.
  \item When \texttt{gatewayUrl} is set, the UI does not fall back to config or environment credentials.
\end{itemize}

Provide \texttt{token} (or \texttt{password}) explicitly. Missing explicit credentials is an error.
\begin{itemize}
  \item Use \texttt{wss://} when the Gateway is behind TLS (Tailscale Serve, HTTPS proxy, etc.).
  \item \texttt{gatewayUrl} is only accepted in a top-level window (not embedded) to prevent clickjacking.
  \item Non-loopback Control UI deployments must set \texttt{gateway.controlUi.allowedOrigins}
\end{itemize}

explicitly (full origins). This includes remote dev setups.
\begin{itemize}
  \item \texttt{gateway.controlUi.dangerouslyAllowHostHeaderOriginFallback=true} enables
\end{itemize}

Host-header origin fallback mode, but it is a dangerous security mode.

Example:

\begin{lstlisting}[language=json]
{
  gateway: {
    controlUi: {
      allowedOrigins: ["http://localhost:5173"],
    },
  },
}
\end{lstlisting}


Remote access setup details: \href{/gateway/remote}{Remote access}.


\clearpage

% === Source file: web/dashboard.md ===
\section{Dashboard (Control UI)}


The Gateway dashboard is the browser Control UI served at \texttt{/} by default
(override with \texttt{gateway.controlUi.basePath}).

Quick open (local Gateway):

\begin{itemize}
  \item \href{http://127.0.0.1:18789/}{http://127.0.0.1:18789/} (or \href{http://localhost:18789/}{http://localhost:18789/})
\end{itemize}


Key references:

\begin{itemize}
  \item \href{/web/control-ui}{Control UI} for usage and UI capabilities.
  \item \href{/gateway/tailscale}{Tailscale} for Serve/Funnel automation.
  \item \href{/web}{Web surfaces} for bind modes and security notes.
\end{itemize}


Authentication is enforced at the WebSocket handshake via \texttt{connect.params.auth}
(token or password). See \texttt{gateway.auth} in \href{/gateway/configuration}{Gateway configuration}.

Security note: the Control UI is an \textbf{admin surface} (chat, config, exec approvals).
Do not expose it publicly. The UI stores the token in \texttt{localStorage} after first load.
Prefer localhost, Tailscale Serve, or an SSH tunnel.

\subsection{Fast path (recommended)}


\begin{itemize}
  \item After onboarding, the CLI auto-opens the dashboard and prints a clean (non-tokenized) link.
  \item Re-open anytime: \texttt{openclaw dashboard} (copies link, opens browser if possible, shows SSH hint if headless).
  \item If the UI prompts for auth, paste the token from \texttt{gateway.auth.token} (or \texttt{OPENCLAW\_GATEWAY\_TOKEN}) into Control UI settings.
\end{itemize}


\subsection{Token basics (local vs remote)}


\begin{itemize}
  \item \textbf{Localhost}: open \texttt{http://127.0.0.1:18789/}.
  \item \textbf{Token source}: \texttt{gateway.auth.token} (or \texttt{OPENCLAW\_GATEWAY\_TOKEN}); the UI stores a copy in localStorage after you connect.
  \item \textbf{Not localhost}: use Tailscale Serve (tokenless for Control UI/WebSocket if \texttt{gateway.auth.allowTailscale: true}, assumes trusted gateway host; HTTP APIs still need token/password), tailnet bind with a token, or an SSH tunnel. See \href{/web}{Web surfaces}.
\end{itemize}


\subsection{If you see “unauthorized” / 1008}


\begin{itemize}
  \item Ensure the gateway is reachable (local: \texttt{openclaw status}; remote: SSH tunnel \texttt{ssh -N -L 18789:127.0.0.1:18789 user@host} then open \texttt{http://127.0.0.1:18789/}).
  \item Retrieve the token from the gateway host: \texttt{openclaw config get gateway.auth.token} (or generate one: \texttt{openclaw doctor --generate-gateway-token}).
  \item In the dashboard settings, paste the token into the auth field, then connect.
\end{itemize}



\clearpage

% === Source file: web/index.md ===
\section{Web (Gateway)}


The Gateway serves a small \textbf{browser Control UI} (Vite + Lit) from the same port as the Gateway WebSocket:

\begin{itemize}
  \item default: \texttt{http://:18789/}
  \item optional prefix: set \texttt{gateway.controlUi.basePath} (e.g. \texttt{/openclaw})
\end{itemize}


Capabilities live in \href{/web/control-ui}{Control UI}.
This page focuses on bind modes, security, and web-facing surfaces.

\subsection{Webhooks}


When \texttt{hooks.enabled=true}, the Gateway also exposes a small webhook endpoint on the same HTTP server.
See \href{/gateway/configuration}{Gateway configuration} → \texttt{hooks} for auth + payloads.

\subsection{Config (default-on)}


The Control UI is \textbf{enabled by default} when assets are present (\texttt{dist/control-ui}).
You can control it via config:

\begin{lstlisting}[language=json]
{
  gateway: {
    controlUi: { enabled: true, basePath: "/openclaw" }, // basePath optional
  },
}
\end{lstlisting}


\subsection{Tailscale access}


\subsubsection{Integrated Serve (recommended)}


Keep the Gateway on loopback and let Tailscale Serve proxy it:

\begin{lstlisting}[language=json]
{
  gateway: {
    bind: "loopback",
    tailscale: { mode: "serve" },
  },
}
\end{lstlisting}


Then start the gateway:

\begin{lstlisting}[language=bash]
openclaw gateway
\end{lstlisting}


Open:

\begin{itemize}
  \item \texttt{https:///} (or your configured \texttt{gateway.controlUi.basePath})
\end{itemize}


\subsubsection{Tailnet bind + token}


\begin{lstlisting}[language=json]
{
  gateway: {
    bind: "tailnet",
    controlUi: { enabled: true },
    auth: { mode: "token", token: "your-token" },
  },
}
\end{lstlisting}


Then start the gateway (token required for non-loopback binds):

\begin{lstlisting}[language=bash]
openclaw gateway
\end{lstlisting}


Open:

\begin{itemize}
  \item \texttt{http://:18789/} (or your configured \texttt{gateway.controlUi.basePath})
\end{itemize}


\subsubsection{Public internet (Funnel)}


\begin{lstlisting}[language=json]
{
  gateway: {
    bind: "loopback",
    tailscale: { mode: "funnel" },
    auth: { mode: "password" }, // or OPENCLAW_GATEWAY_PASSWORD
  },
}
\end{lstlisting}


\subsection{Security notes}


\begin{itemize}
  \item Gateway auth is required by default (token/password or Tailscale identity headers).
  \item Non-loopback binds still \textbf{require} a shared token/password (\texttt{gateway.auth} or env).
  \item The wizard generates a gateway token by default (even on loopback).
  \item The UI sends \texttt{connect.params.auth.token} or \texttt{connect.params.auth.password}.
  \item For non-loopback Control UI deployments, set \texttt{gateway.controlUi.allowedOrigins}
\end{itemize}

explicitly (full origins). Without it, gateway startup is refused by default.
\begin{itemize}
  \item \texttt{gateway.controlUi.dangerouslyAllowHostHeaderOriginFallback=true} enables
\end{itemize}

Host-header origin fallback mode, but is a dangerous security downgrade.
\begin{itemize}
  \item With Serve, Tailscale identity headers can satisfy Control UI/WebSocket auth
\end{itemize}

when \texttt{gateway.auth.allowTailscale} is \texttt{true} (no token/password required).
HTTP API endpoints still require token/password. Set
\texttt{gateway.auth.allowTailscale: false} to require explicit credentials. See
\href{/gateway/tailscale}{Tailscale} and \href{/gateway/security}{Security}. This
tokenless flow assumes the gateway host is trusted.
\begin{itemize}
  \item \texttt{gateway.tailscale.mode: "funnel"} requires \texttt{gateway.auth.mode: "password"} (shared password).
\end{itemize}


\subsection{Building the UI}


The Gateway serves static files from \texttt{dist/control-ui}. Build them with:

\begin{lstlisting}[language=bash]
pnpm ui:build # auto-installs UI deps on first run
\end{lstlisting}



\clearpage

% === Source file: web/tui.md ===
\section{TUI (Terminal UI)}


\subsection{Quick start}


\begin{enumerate}
  \item Start the Gateway.
\end{enumerate}


\begin{lstlisting}[language=bash]
openclaw gateway
\end{lstlisting}


\begin{enumerate}
  \item Open the TUI.
\end{enumerate}


\begin{lstlisting}[language=bash]
openclaw tui
\end{lstlisting}


\begin{enumerate}
  \item Type a message and press Enter.
\end{enumerate}


Remote Gateway:

\begin{lstlisting}[language=bash]
openclaw tui --url ws://<host>:<port> --token <gateway-token>
\end{lstlisting}


Use \texttt{--password} if your Gateway uses password auth.

\subsection{What you see}


\begin{itemize}
  \item Header: connection URL, current agent, current session.
  \item Chat log: user messages, assistant replies, system notices, tool cards.
  \item Status line: connection/run state (connecting, running, streaming, idle, error).
  \item Footer: connection state + agent + session + model + think/verbose/reasoning + token counts + deliver.
  \item Input: text editor with autocomplete.
\end{itemize}


\subsection{Mental model: agents + sessions}


\begin{itemize}
  \item Agents are unique slugs (e.g. \texttt{main}, \texttt{research}). The Gateway exposes the list.
  \item Sessions belong to the current agent.
  \item Session keys are stored as \texttt{agent::}.
  \item If you type \texttt{/session main}, the TUI expands it to \texttt{agent::main}.
  \item If you type \texttt{/session agent:other:main}, you switch to that agent session explicitly.
  \item Session scope:
  \item \texttt{per-sender} (default): each agent has many sessions.
  \item \texttt{global}: the TUI always uses the \texttt{global} session (the picker may be empty).
  \item The current agent + session are always visible in the footer.
\end{itemize}


\subsection{Sending + delivery}


\begin{itemize}
  \item Messages are sent to the Gateway; delivery to providers is off by default.
  \item Turn delivery on:
  \item \texttt{/deliver on}
  \item or the Settings panel
  \item or start with \texttt{openclaw tui --deliver}
\end{itemize}


\subsection{Pickers + overlays}


\begin{itemize}
  \item Model picker: list available models and set the session override.
  \item Agent picker: choose a different agent.
  \item Session picker: shows only sessions for the current agent.
  \item Settings: toggle deliver, tool output expansion, and thinking visibility.
\end{itemize}


\subsection{Keyboard shortcuts}


\begin{itemize}
  \item Enter: send message
  \item Esc: abort active run
  \item Ctrl+C: clear input (press twice to exit)
  \item Ctrl+D: exit
  \item Ctrl+L: model picker
  \item Ctrl+G: agent picker
  \item Ctrl+P: session picker
  \item Ctrl+O: toggle tool output expansion
  \item Ctrl+T: toggle thinking visibility (reloads history)
\end{itemize}


\subsection{Slash commands}


Core:

\begin{itemize}
  \item \texttt{/help}
  \item \texttt{/status}
  \item \texttt{/agent } (or \texttt{/agents})
  \item \texttt{/session } (or \texttt{/sessions})
  \item \texttt{/model } (or \texttt{/models})
\end{itemize}


Session controls:

\begin{itemize}
  \item \texttt{/think }
  \item \texttt{/verbose }
  \item \texttt{/reasoning }
  \item \texttt{/usage }
  \item \texttt{/elevated } (alias: \texttt{/elev})
  \item \texttt{/activation }
  \item \texttt{/deliver }
\end{itemize}


Session lifecycle:

\begin{itemize}
  \item \texttt{/new} or \texttt{/reset} (reset the session)
  \item \texttt{/abort} (abort the active run)
  \item \texttt{/settings}
  \item \texttt{/exit}
\end{itemize}


Other Gateway slash commands (for example, \texttt{/context}) are forwarded to the Gateway and shown as system output. See \href{/tools/slash-commands}{Slash commands}.

\subsection{Local shell commands}


\begin{itemize}
  \item Prefix a line with \texttt{!} to run a local shell command on the TUI host.
  \item The TUI prompts once per session to allow local execution; declining keeps \texttt{!} disabled for the session.
  \item Commands run in a fresh, non-interactive shell in the TUI working directory (no persistent \texttt{cd}/env).
  \item Local shell commands receive \texttt{OPENCLAW\_SHELL=tui-local} in their environment.
  \item A lone \texttt{!} is sent as a normal message; leading spaces do not trigger local exec.
\end{itemize}


\subsection{Tool output}


\begin{itemize}
  \item Tool calls show as cards with args + results.
  \item Ctrl+O toggles between collapsed/expanded views.
  \item While tools run, partial updates stream into the same card.
\end{itemize}


\subsection{History + streaming}


\begin{itemize}
  \item On connect, the TUI loads the latest history (default 200 messages).
  \item Streaming responses update in place until finalized.
  \item The TUI also listens to agent tool events for richer tool cards.
\end{itemize}


\subsection{Connection details}


\begin{itemize}
  \item The TUI registers with the Gateway as \texttt{mode: "tui"}.
  \item Reconnects show a system message; event gaps are surfaced in the log.
\end{itemize}


\subsection{Options}


\begin{itemize}
  \item \texttt{--url }: Gateway WebSocket URL (defaults to config or \texttt{ws://127.0.0.1:})
  \item \texttt{--token }: Gateway token (if required)
  \item \texttt{--password }: Gateway password (if required)
  \item \texttt{--session }: Session key (default: \texttt{main}, or \texttt{global} when scope is global)
  \item \texttt{--deliver}: Deliver assistant replies to the provider (default off)
  \item \texttt{--thinking }: Override thinking level for sends
  \item \texttt{--timeout-ms }: Agent timeout in ms (defaults to \texttt{agents.defaults.timeoutSeconds})
\end{itemize}


Note: when you set \texttt{--url}, the TUI does not fall back to config or environment credentials.
Pass \texttt{--token} or \texttt{--password} explicitly. Missing explicit credentials is an error.

\subsection{Troubleshooting}


No output after sending a message:

\begin{itemize}
  \item Run \texttt{/status} in the TUI to confirm the Gateway is connected and idle/busy.
  \item Check the Gateway logs: \texttt{openclaw logs --follow}.
  \item Confirm the agent can run: \texttt{openclaw status} and \texttt{openclaw models status}.
  \item If you expect messages in a chat channel, enable delivery (\texttt{/deliver on} or \texttt{--deliver}).
  \item \texttt{--history-limit }: History entries to load (default 200)
\end{itemize}


\subsection{Connection troubleshooting}


\begin{itemize}
  \item \texttt{disconnected}: ensure the Gateway is running and your \texttt{--url/--token/--password} are correct.
  \item No agents in picker: check \texttt{openclaw agents list} and your routing config.
  \item Empty session picker: you might be in global scope or have no sessions yet.
\end{itemize}



\clearpage

% === Source file: web/webchat.md ===
\section{WebChat (Gateway WebSocket UI)}


Status: the macOS/iOS SwiftUI chat UI talks directly to the Gateway WebSocket.

\subsection{What it is}


\begin{itemize}
  \item A native chat UI for the gateway (no embedded browser and no local static server).
  \item Uses the same sessions and routing rules as other channels.
  \item Deterministic routing: replies always go back to WebChat.
\end{itemize}


\subsection{Quick start}


\begin{enumerate}
  \item Start the gateway.
  \item Open the WebChat UI (macOS/iOS app) or the Control UI chat tab.
  \item Ensure gateway auth is configured (required by default, even on loopback).
\end{enumerate}


\subsection{How it works (behavior)}


\begin{itemize}
  \item The UI connects to the Gateway WebSocket and uses \texttt{chat.history}, \texttt{chat.send}, and \texttt{chat.inject}.
  \item \texttt{chat.history} is bounded for stability: Gateway may truncate long text fields, omit heavy metadata, and replace oversized entries with \texttt{[chat.history omitted: message too large]}.
  \item \texttt{chat.inject} appends an assistant note directly to the transcript and broadcasts it to the UI (no agent run).
  \item Aborted runs can keep partial assistant output visible in the UI.
  \item Gateway persists aborted partial assistant text into transcript history when buffered output exists, and marks those entries with abort metadata.
  \item History is always fetched from the gateway (no local file watching).
  \item If the gateway is unreachable, WebChat is read-only.
\end{itemize}


\subsection{Control UI agents tools panel}


\begin{itemize}
  \item The Control UI \texttt{/agents} Tools panel fetches a runtime catalog via \texttt{tools.catalog} and labels each
\end{itemize}

tool as \texttt{core} or \texttt{plugin:} (plus \texttt{optional} for optional plugin tools).
\begin{itemize}
  \item If \texttt{tools.catalog} is unavailable, the panel falls back to a built-in static list.
  \item The panel edits profile and override config, but effective runtime access still follows policy
\end{itemize}

precedence (\texttt{allow}/\texttt{deny}, per-agent and provider/channel overrides).

\subsection{Remote use}


\begin{itemize}
  \item Remote mode tunnels the gateway WebSocket over SSH/Tailscale.
  \item You do not need to run a separate WebChat server.
\end{itemize}


\subsection{Configuration reference (WebChat)}


Full configuration: \href{/gateway/configuration}{Configuration}

Channel options:

\begin{itemize}
  \item No dedicated \texttt{webchat.*} block. WebChat uses the gateway endpoint + auth settings below.
\end{itemize}


Related global options:

\begin{itemize}
  \item \texttt{gateway.port}, \texttt{gateway.bind}: WebSocket host/port.
  \item \texttt{gateway.auth.mode}, \texttt{gateway.auth.token}, \texttt{gateway.auth.password}: WebSocket auth (token/password).
  \item \texttt{gateway.auth.mode: "trusted-proxy"}: reverse-proxy auth for browser clients (see \href{/gateway/trusted-proxy-auth}{Trusted Proxy Auth}).
  \item \texttt{gateway.remote.url}, \texttt{gateway.remote.token}, \texttt{gateway.remote.password}: remote gateway target.
  \item \texttt{session.*}: session storage and main key defaults.
\end{itemize}



\clearpage
